\documentclass[12pt]{article}
\usepackage{amsmath,amsthm,amssymb}
\usepackage{enumitem}
\title{Abstract Algebra - Homework 2}
\author{Simon Gustafsson}
\date{}

\begin{document}
\maketitle

\section*{Problem 1}
\begin{enumerate}[label=(\arabic*)]
    \item    
    \[
    ij = k, \quad ij = -ji \implies -ji = k \implies ji = -k
    \]
    Hence, \( ij \neq ji \), so multiplication is not commutative and \( \mathbb{H} \) is a non-commutative ring.

    \item
    \textbf{(a)} Let \( x, y \in \mathbb{H} \). One can verify by direct computation that conjugation reverses the order of multiplication,
    \( \overline{xy} = \bar{y} \bar{x}.\)
    Note that for any \( h = a + bi + cj + dk \in \mathbb{H} \), we have
    \[
    N(h) = h\bar{h} = a^2 + b^2 + c^2 + d^2 \in \mathbb{R},
    \]
    so \( N(h) \) lies in the center of \( \mathbb{H} \) and commutes with all elements. Using this, we compute:
    \[
    \begin{aligned}
        N(xy) &= (xy)\overline{(xy)} \\
              &= (xy)(\bar{y}\bar{x}) \\
              &= x(y\bar{y})\bar{x} \\
              &= xN(y)\bar{x} \\
              &= N(y) x\bar{x} \\
              &= N(y)N(x) = N(x)N(y).
    \end{aligned}
    \]

    \textbf{(b)} Since \( N(x) = x\bar{x} = a^2 + b^2 + c^2 + d^2 \), it follows that
    \[
    N(x) = 0 \iff a = b = c = d = 0 \iff x = 0.
    \]
    Thus, \( N(x) = 0 \) if and only if \( x = 0 \).
    \item
    Let \( x \in \mathbb{H} \) be nonzero. From previous section, we have \(N(x) = x\bar{x} = a^2 + b^2 + c^2 + d^2 > 0.\)
    Define the inverse of \( x \) by
    \[
    x^{-1} := \frac{\bar{x}}{N(x)}.
    \]
    Then:
    \[
    x x^{-1} = x \cdot \frac{\bar{x}}{N(x)} = \frac{x\bar{x}}{N(x)} = \frac{N(x)}{N(x)} = 1,
    \quad
    x^{-1} x = \frac{\bar{x}}{N(x)} \cdot x = \frac{\bar{x}x}{N(x)} = \frac{N(x)}{N(x)} = 1.
    \]
    Hence, every nonzero element of \( \mathbb{H} \) is invertible.

    Suppose \( x, y \in \mathbb{H} \) with \( x \ne 0 \) and \( xy = 0 \). Then, \(x^{-1}(xy) = (x^{-1}x)y = 1 \cdot y = y = 0.\)
    Thus, \( y = 0 \), and similarly one shows that \( x = 0 \) if \( y \ne 0 \) and \( xy = 0 \). Therefore, \( \mathbb{H} \) has no zero-divisors, and every nonzero element is invertible. 
\end{enumerate}



\end{document}